\section{Datasets and OOD Evaluation Protocols}
\label{sec:datasets}

The experiments use converted OpenEW-SA artifacts with one metadata row per sample, an index into a shared feature array, and symbolic labels loaded as strings. Each frozen protocol contains an ID training split, an ID validation split, an ID test split, and an OOD test split. The evaluation CSV is the concatenation of test ID and test OOD rows only. Table~\ref{tab:dataset-protocols} reports the verified evaluation counts used by every compared v3 method.

\begin{table*}[!t]
  \caption{Dataset and OOD protocol summary. Processed input view, recognition target, frozen OOD definition, and verified ID/OOD evaluation counts for ElectroSense, DeepSense, and JamShield. Counts refer to the frozen evaluation set, not the training or validation partitions.}
  \label{tab:dataset-protocols}
  \centering
  \scriptsize
  \setlength{\tabcolsep}{3pt}
  \renewcommand{\arraystretch}{1.12}
  \begin{tabularx}{\textwidth}{@{}l l l l Y Y r r r@{}}
    \toprule
    Dataset & Input & Target & Protocol & ID definition & OOD definition & ID & OOD & Total \\
    \midrule
    ElectroSense & PSD & Signal technology & Class OOD &
      DAB, DVB-T, FM, and LTE classes &
      GSM and TETRA classes &
      5,840 & 16,550 & 22,390 \\
    DeepSense & I/Q & Wi-Fi occupancy code & Acquisition-day OOD &
      Retained day-one samples &
      Held-out day-two samples &
      3,200 & 16,000 & 19,200 \\
    JamShield & Tabular radio/network metrics & Normal versus abnormal interference & Scenario OOD &
      Retained benign and jammer scenarios &
      Held-out benign and jammer scenarios &
      14,534 & 19,817 & 34,351 \\
    \bottomrule
  \end{tabularx}
\end{table*}


\subsection{ElectroSense Class OOD}

ElectroSense provides crowdsensed PSD data and an official wireless-technology classification framework~\cite{rajendran2018electrosense,scalingi2023framework,scalingi2023electrosensepsd}. The taxonomy and split used here are scoped to the converted OpenEW-SA subset: DAB, DVB-T, FM, and LTE are treated as known classes for classifier fitting and ID evaluation, while GSM and TETRA are withheld from classifier training and treated as OOD at evaluation. The upstream framework code also defines an \method{unkn} label; the analyzed frozen manifest and split artifacts contain only the six named technology labels and therefore exclude \method{unkn}. This six-class construction is not attributed to the Zenodo record itself.

\subsection{DeepSense Acquisition-Day OOD}

DeepSense contributes processed I/Q windows labeled by binary occupancy codes~\cite{uvaydov2021deepsense,wineslab2021deepsensedataset}. All occupancy codes remain known labels, while the acquisition day defines the shift: retained day-one rows provide train, validation, and test-ID samples, and held-out day-two rows provide OOD samples. Labels such as \method{0000}, \method{0001}, \method{0010}, and \method{0100} are identifiers and are read and written as strings so their leading zeros cannot be lost. This protocol tests domain shift rather than unseen-class rejection; the official dataset record documents different transmitter orientations across the two acquisition days.

\subsection{JamShield Scenario OOD}

JamShield contributes tabular radio and network telemetry for normal-versus-abnormal interference recognition~\cite{panitsas2025jamshield,panitsas2024jamshielddataset}. Both recognition labels occur in the retained scenarios, while a frozen subset of benign and jammer scenarios is held out as OOD. The task therefore asks whether a classifier and its training geometry can identify scenario novelty even when the nominal label vocabulary remains unchanged. The exact retained and held-out domain identifiers are fixed in the split manifests. The peer-reviewed ICC 2025 publication is cited for the system, while the public dataset record supports the scenario and telemetry provenance; no explicit dataset license was identified in the verified public records.

\subsection{Leakage Controls}

Training features and labels are used to fit the classifier and the class centroids. ID validation predictions are used to fit temperature scaling, and ID validation component scores are used to fit score normalization. Test-ID and test-OOD labels are used only for final metric computation, audit alignment, and the explicitly labeled post-hoc DeepSense diagnostic. They are not used to choose the score orientation, normalization statistics, fusion weights, primary method, comparator set, or a deployment threshold.
